\section{Limitations}\label{sec:7}\label{limitations}

The theory links decision gap, coordinate edit capacity and necessary cost through an elementary coupling argument. The lower and upper bounds are not matching and leave the variance-sensitive intermediate regime unresolved. Consequently, observed near-census stopping times are not established to be optimal. Side information that rules out the constructed alternatives can also remove the stated obstruction.

The public matrices contain upstream filtering, older model families and dependent model lineages. Margin-nearest selection deliberately enriches difficult pairs and cannot estimate their prevalence. We use item-weighted frozen scores, not official macro-weighted leaderboard targets, future generations or new-task performance. Unequal-cost evidence is synthetic; no dollar savings are established.

Our strongest empirical comparisons are common-target replays of explicit inference constructions. Arviv's efficacy target differs from our practical-superiority target, so an original-pipeline comparison requires an explicit adaptation and calibration check. Predictors and strata in Appendix~\ref{app:c} are adaptations, not an exhaustive comparison against every configuration of CELEUS or PULSE. No state-of-the-art cost claim follows. This missing harmonized comparison limits conclusions about relative efficiency, even though it does not affect the validity of the inclusion bound.
